%  ============================================================================
% MAIN DOCUMENT: Statistical Mechanics of Money
% A paper on applying Boltzmann-Gibbs distributions to economic systems
% ============================================================================

\documentclass[10pt, a4paper, twocolumn]{article}

% Include shared configuration
% ============================================================================
% PREAMBLE: Shared configuration and packages for the paper
% ============================================================================

% ---- DOCUMENT PACKAGES ----
\usepackage[utf8]{inputenc}
\usepackage[T1]{fontenc}
\usepackage{lmodern}
\usepackage[a4paper, margin=1in]{geometry}

% ---- MATHEMATICS & PHYSICS ----
\usepackage{amsmath}
\usepackage{amssymb}
\usepackage{physics}  % For \dv, \pdv, \dd, etc.
\usepackage{siunitx}  % For units formatting

% ---- GRAPHICS & FIGURES ----
\usepackage{graphicx}
\usepackage{tikz}
\usepackage{pgfplots}
\pgfplotsset{compat=1.17}
\graphicspath{{../figures/}}

% ---- CITATIONS & BIBLIOGRAPHY ----
\usepackage[square, numbers, sort&compress]{natbib}
\bibliographystyle{plainnat}

% ---- HYPERLINKS & CROSS-REFERENCES ----
\usepackage{hyperref}
\hypersetup{
    colorlinks=true,
    linkcolor=blue,
    citecolor=blue,
    urlcolor=blue,
    bookmarksnumbered=true
}

% ---- CODE LISTINGS ----
\usepackage{listings}
\usepackage{xcolor}

\lstset{
    language=Python,
    basicstyle=\ttfamily\small,
    keywordstyle=\color{blue}\bfseries,
    commentstyle=\color{gray}\itshape,
    stringstyle=\color{red},
    backgroundcolor=\color{gray!10},
    breaklines=true,
    breakatwhitespace=true,
    showstringspaces=false,
    numbers=left,
    numberstyle=\tiny\color{gray},
    numbersep=5pt,
    frame=single,
    framerule=0.5pt,
    captionpos=b,
    tabsize=4
}

% ---- THEOREMS & PROOF ENVIRONMENTS ----
\usepackage{amsthm}
\theoremstyle{definition}
\newtheorem{definition}{Definition}[section]
\newtheorem{theorem}{Theorem}[section]
\newtheorem{lemma}[theorem]{Lemma}
\newtheorem{corollary}[theorem]{Corollary}
\newtheorem{proposition}[theorem]{Proposition}

\theoremstyle{remark}
\newtheorem{remark}{Remark}[section]
\newtheorem{example}{Example}[section]

% ---- PHYSICS NOTATION (Custom Commands) ----
\newcommand{\Temperature}{T}
\newcommand{\Entropy}{S}
\newcommand{\Money}{m}
\newcommand{\Agents}{N}
\newcommand{\TotalMoney}{M}
\newcommand{\AvgMoney}{\langle m \rangle}
\newcommand{\Prob}[1]{P\left(#1\right)}
\newcommand{\Expected}[1]{\langle #1 \rangle}
\newcommand{\MarkovChain}{\{X_n\}_{n \geq 0}}
\newcommand{\GibbsEntropy}{S = -\sum_i p_i \ln p_i}
\newcommand{\Boltzmann}{k_B}
\newcommand{\Partition}{Z}

% ---- DOCUMENT FORMATTING ----
\usepackage{fancyhdr}
\usepackage{setspace}
\onehalfspacing

\pagestyle{fancy}
\fancyhf{}
\fancyhead[L]{\leftmark}
\fancyhead[R]{\thepage}
\renewcommand{\headrulewidth}{0.4pt}

% ---- CAPTIONS & FLOATS ----
\usepackage{caption}
\captionsetup{
    font=small,
    labelfont=bf,
    justification=centering
}

% ---- APPENDIX ----
\usepackage{appendix}

% ============================================================================
% END OF PREAMBLE
% ============================================================================


% ============================================================================
% TITLE PAGE
% ============================================================================
\title{
    \textbf{Statistical Mechanics of Money:} \\
    \large Applying Boltzmann-Gibbs Distributions and Entropy \\
    to Economic Dynamics via Agent-Based Simulation
}

\author{
    \textit{Submitted for:} \\
    Physics of Complex Systems (Exam) \\
    Università della Magistrale in Fisica
}

\date{\today}

% ============================================================================
% DOCUMENT BEGINS
% ============================================================================
\begin{document}

\maketitle

% ---- TABLE OF CONTENTS ----
\tableofcontents
\newpage

% ============================================================================
% CHAPTERS
% ============================================================================

% Chapter 0: Abstract (Full Width)
%  ============================================================================
% ABSTRACT
% ============================================================================

\section*{Abstract}


\noindent \textbf{Keywords:} Statistical mechanics, Boltzmann-Gibbs distribution, Gibbs entropy, Agent-based simulation, Markov processes, Economic dynamics, Complex systems.

\newpage


% The document class already uses a two-column layout for the rest of the paper.
% Chapter 1: Introduction
%  ============================================================================
% CHAPTER 1: INTRODUCTION
% ============================================================================

\section{Introduction}
This paper investigates the applicability of \textbf{statistical mechanics} to economic systems, specifically how money distribution in a society of interacting economic agents follows the \textbf{Boltzmann-Gibbs distribution}. 
This is a foundational principle from thermodynamics and statistical physics. 
The assumption is that money, like energy in a physical system, is conserved during transactions between agents.

I conducted agent based simulations modeling economic transactions as random exchanges of money between agents, analogous to particle collisions in a gas. Under various transaction mechanisms (constant exchanges, fraction of individual wealth, fraction of system average), the simulated money distribution converges to an exponential form $\Prob{m} = \frac{1}{\Temperature} e^{-m/\Temperature}$, where $\Temperature = M/N$ represents the system's effective ``temperature'' (average wealth per agent).

Three simulation models are analyzed:
\begin{enumerate}
    \item \textbf{Basic money distribution} with no debt constraints
    \item \textbf{System with debt} allowing limited negative balances
%    \item \textbf{Stock market extension} incorporating both cash and equity trades with price dynamics
\end{enumerate}

We demonstrate that \textbf{Gibbs entropy} correctly predicts equilibrium states and that the system evolves as a \textbf{Markov process} toward maximum entropy. 
These results validate the hypothesis that economic inequality and wealth concentration emerge as natural equilibrium
states when agents interact via random transactions, analogously to how physical systems spontaneously approach thermodynamic equilibrium.

This paper connects fundamental concepts from statistical mechanics from Markov processes, Boltzmann-Gibbs distributions, Gibbs entropy, and dynamical systems theory
to economics and the so called "econophysics" or "social physics", providing a computational validation of theoretical predictions.

\vspace{1em}


\subsection{Motivation}



\newpage


% Chapter 2: Model & Theory
%  ============================================================================
% CHAPTER 2: MODEL AND THEORY
% ============================================================================

\section{Model and Theory}
As previously discussed, we can assume that agents can be modeled as particles in a closed system, where money is conserved.
Instead of exchanging energy, agents exchange money through random transactions. 
The system evolves through random transactions, forming a Markov process. 
The theoretical framework predicts that the wealth distribution converges to the Boltzmann-Gibbs distribution, with entropy increasing over time until equilibrium is reached.

\subsection{Conservation of Money}
One could assume that Boltzmann's derivation of the energy distribution can be applied to other statistical systems with a conserved quantity, such as money.
Indeed, we can consider the economy as a closed system with a fixed number of agents $\Agents$ and a total amount of money $\TotalMoney$.
Therefore, the total money in the system is conserved.
As a starting point, we can consider a simple model where agents exchange money through random transactions, and the total amount of money remains constant.

Introducing the concept of money conservation we can assume the model \textbf{without debt}, where agents cannot have negative balances.
So the conditions for this model are:
\begin{itemize}
    \item $\displaystyle \sum_{i=1}^{\Agents} m_i = \TotalMoney$.
    \item $m_i \geq 0$, for $i = 1, \ldots, \Agents$.
    \item $\delta m_i = -\delta m_j$, for $i \neq j$.
    \item $\delta_m < m_i$.
\end{itemize}

\subsection{Boltzmann-Gibbs Distribution of Money}
As argued in previous section, we can associate the money distribution
to the energy exchange in a physical system. 
For this reason, we can assume that the probability of finding an agent with money $m$ is given by the Boltzmann-Gibbs distribution:
\begin{equation}
    \Prob{m} = c e^{-m/\Temperature},
\end{equation}
where $c$ is a normalization constant and $\Temperature$ is the effective temperature of the system: $T = \TotalMoney / N$, which is average money per agent.

The \textbf{entropy} of the system is defined as:
\begin{equation}
    S = -\sum_{i=1}^{\Agents} \Prob{m_i} \log \Prob{m_i}.
\end{equation}
The \textbf{equilibrium} is reached when the entropy is maximized when the system "cools off" to the Boltzmann-Gibbs distribution.

\subsection{Transaction Mechanisms}
Many different kinds of transactions can be considered, such as:
\begin{itemize}
    \item \textbf{Constant transactions: $\Delta m = 1$} where a fixed amount of money is exchanged between agents. 
    \item \textbf{Fraction of Pair's Average:} A random fraction of the average money of the pair $\frac{m_i + m_j}{2}$ is exchanged:
                                                $\Delta m = r \cdot \frac{m_i + m_j}{2}$, 
    \item \textbf{Fractional transactions} where a fraction $\gamma$ of the agent's money is exchanged, the two agents will exchange
                                            only a portion of their money. This mechanism show how richer agents can exchange more money than poorer agents, and it is more realistic.
\end{itemize}

\subsection{The different models}
Different assumptions were made in the literature to model the money exchange between agents.
The aim it's to find a model that can reproduce the empirical wealth distribution observed in real economies.
\subsubsection{Model without debt}
The simplest model has constant transactions and no debt can be accumulated. This model leads to the Boltzmann-Gibbs distribution as previously discussed.
\subsubsection{Model with debt}
The exact mechanisms of limiting debt in the real economy are complicated and obscured.
In this scenario, debt is introduced in the model \cite{yakovenko2012statistical}. 




Therefore we can add a saving propensity $\lambda$ to the model, where agents save a fraction of their money before exchanging it.


% Chapter 3: Results
%  ============================================================================
% CHAPTER 3: RESULTS
% ============================================================================

\section{Results}

\subsection{Basic Money Distribution}

In this scenario, we consider a simple model of money distribution among agents without debt. 
The total money of the system is conserved, and no agent can have a negative balance.
The agents engage with each other randomly but always with a fixed amount of money in each transaction.

Said so, we can assume that the money distribution will converge to the Boltzmann-Gibbs distribution as predicted by the theoretical framework.
It is possible to see the results in Figure \ref{fig:NoDebtImage}, where the histogram of the money distribution is plotted along with the theoretical Boltzmann-Gibbs distribution.

\begin{figure}[H]
    \centering
    \includegraphics[width=0.4\textwidth]{NoDebtImage.png}
    \caption{The histogram shows the money distribution of agents after a large number of transactions, while the red line represents the theoretical Boltzmann-Gibbs distribution.
                The simulation was performed with $N=3500$, $M = 10 000$, and $\Delta m = 1$. The simulation was run for $10^6$ transactions.}
    \label{fig:NoDebtImage}
\end{figure}

Later on it will be discussed how the richer agents don't really follow Boltzmann-Gibbs distribution, but they follow a power law distribution, as observed in real economies.
These agents are not in equilibrium with the rest of the system, therefore they are not part of the thermalized system.

\subsection{Money Distribution with Debt}


\subsection{Power Law Distribution}




% Chapter 4: Conclusion
%  ============================================================================
% CHAPTER 4: CONCLUSION
% ============================================================================

\section{Conclusion}

\subsection{Summary}

[Conclusion content to be filled...]

\newpage


% ============================================================================
% APPENDICES
% ============================================================================
\appendix

\section{Mathematical Details}\label{app:mathematics}
%  ============================================================================
% APPENDIX A: MATHEMATICAL DETAILS
% ============================================================================

\subsection{Boltzmann-Gibbs Distribution}

[Mathematical details to be filled...]


\subsection{Reverse Ratio choice of value $r$}
From section \ref{sec:debtModel}, we can see that the choice of $r$ is crucial for the model to work properly.
The choice of $r$ is based on the assumption that the total debt in the system is limited by the required reserve ratio.
To have matching populations at $m=0$, the condition $P(m_+)=P(m_-)$ reduces to
\[
\frac{N_+}{T_+} = \frac{N_-}{T_-},
\]
with
\[
T_+ = \frac{M_+}{N_+}, \qquad T_- = \frac{M_-}{N_-}.
\]
Using the definitions of the two sectors,
\[
N_+ = rN, \qquad N_- = (1-r)N,
\]
\[
M_+ = \frac{M_0}{r}, \qquad M_- = \frac{M_0(1-r)}{r},
\]
we obtain
\[
T_+ = \frac{M_0}{r^2 N}, \qquad T_- = \frac{M_0}{rN}.
\]
Substituting these expressions into the matching condition yields
\[
r \cdot \frac{r^2N}{M_0} = (1-r)\cdot \frac{rN}{M_0},
\]
which simplifies to
\[
r^2 + r - 1 = 0.
\]
Therefore the reverse ratio that makes the two distributions meet at the origin is
\[
r = \frac{\sqrt{5}-1}{2} \approx 0.618.
\]
This is the value used in the simulations to compare the positive and negative branches of the double-sided Boltzmann-Gibbs distribution.




\section{Code Snippets}\label{app:code}
%  ============================================================================
% APPENDIX B: PYTHON CODE SNIPPETS
% ============================================================================

\subsection{Core Simulation}

[Code snippets to be filled...]


% ============================================================================
% BIBLIOGRAPHY
% ============================================================================
\newpage
\bibliography{bib/references}

% ============================================================================
% END OF DOCUMENT
% ============================================================================
\end{document}

